% Бакалаврын судалгааны тайлан: speaker recognition ба diarization
% Хэрэгжүүлэлт: Монгол улсын их сургуулийн ажилтан оюутны өгөгдөл дээр туршилтууд,
% үр дүнгийн шинжилгээ, дараагийн судалгааны чиглэлүүд.

\section{Удиртгал}
Энэ бүлэгт судалгааны сэдэв, асуудлын хүрээ, зорилго ба ач холбогдолыг танилцуулна. Speaker recognition ба speaker diarization нь одоохи дууны боловсролын чиглэлд чухал асуудал болж байна. Энэхүү судалгааны ажил нь эдгээр асуудлыг Монгол хэлний өгөгдөлд хэрэгжүүлж, нэгэн энгийн, төсөөлөлд тохирсон багц архитектурыг санал болгож байна. 

\section{Асуудлын тодорхойлолт ба хандлага}
Сонирхох асуудал:
- Speaker recognition системийн нарийвчлал, үнэлгээний стратеги;
- Speaker diarization-ийн accuracy, separation capability, real-time applicability.
- Монгол хэлээр контекстуал өгөгдөл дээр ажиллах боломж;
- Өгөгдөл боловсруулалтын шалгуур ба үр дүнгийн үнэлгээ.

\section{Метод ба архитектур}
Энэхүү судалгаанд дараах үндсэн хэсгүүд орно:
- Үгсийн төлөв байдлын төлөөлөл: x-vectors (DNN-т суурилсан эмбеддинг), PLDA зэрэг
- Ярианы хэсгүүдийн ялгаварлал: Agglomerative Hierarchical Clustering (AHC) ба VBx аргачлалын хослол
- Дискриптив preprocessing: Voice Activity Detection, MFCC, features normalization
- Диаризация: сегментчилэлт ба хэлээр ялгах үр дүнгийн ба хэлэлцүүлэл

\section{Өгөгдөл ба үнэлгээ}
Өгөгдөл:
- VoxCeleb1/2 зэрэг олон улсын нээлттэй өгөгдлүүдийг ашиглах буюу xmlns-стандарт өгөгдөл дээр туршилагдах.
- Монгол хэлний өгөгдөл бүрдүүлэх шаардлага байдлаас шалтгаалан нэмэлт өгөгдөл үүсгэх төлөвлөгөө байна.
- Үнэлгээ: DER ба EER зэрэг метрүүд; базелайн системтэй харьцуулалт.

\section{Үр дүн ба хэлэлцүүлэг}
Дасгалжуулалт ба үнэлгээний үр дүнгүүдийг хүснэгт болон зургийн хэлбэрээр танилцуулна. Дүгнэлтэд системийн давуу тал, хязгаарлалтууд ба Монгол хэлэнд зориулсан боломжит өргөтгөлүүдийг авч үзнэ.

\section{Дүгнэлт ба ирээдүйн ажил}
Энэ судалгааны эцсийн дүгнэлт нь таны төсөлд захиралтай, үр дүнгийн талаарх үндсэн ойлголтуудыг өгч, ирээдүйн судалгааны чиглэлүүдийг тодорхойлно.

\section{Ишлэл}
Ашигласан материал, ном зүйг энд оруулна.
